\section{Глава 10. Управление с единой стратегией и аппроксимацией}

\subsection{Упражнение 10.1(с. 294)}
В этой главе мы не рассматривали явно и не приводили псевдокод для методов
Монте-Карло. На что они были бы похожи? Почему не имеет смысла приводить для них псевдокод? Какие результаты они дали бы в задаче о машине на горе?

\begin{solution}
    На что похожи. Метод Монте-Карло --- это частный случай уже приведённого
    $n$-шагового полуградиентного Sarsa, когда окно дотягивается до конца
    эпизода. При $t+n\ge T$ бутстрэп-член в (10.4) исчезает и цель становится
    полным доходом: $G_{t:t+n}=G_t$. Обновление превращается в
    \begin{equation*}
        \mathbf w \leftarrow \mathbf w + \alpha\big[G_t - \hat q(S_t, A_t, \mathbf w)\big]\nabla\hat q(S_t, A_t, \mathbf w),
    \end{equation*}
    применяемое ко всем парам эпизода после его завершения.

    Почему не нужен отдельный псевдокод. Именно потому, что MC не требует ничего
    нового: он буквально получается из $n$-шагового Sarsa, если взять $n$
    настолько большим, что цель всегда становится полным доходом $G_t$. Отдельная
    запись лишь повторила бы уже имеющуюся.

    Результаты на машине на горе. Здесь MC проявил бы себя плохо, и причина
    глубже, чем медленная сходимость. MC обновляет ценности только по завершении
    эпизода --- до этого $G_t$ не определён. Значит весь первый эпизод агент
    блуждает с замороженной начальной политикой и не делает ни одного обновления.
    А при оптимистичном старте ($\mathbf w=0$) и слабой стартовой политике нет
    гарантии, что эпизод завершится в разумное время; в худшем случае MC не
    сделает ни одного обновления вовсе.

    Sarsa этой ловушки лишён: он обновляет $\hat q$ прямо во время эпизода, и эти
    обновления помогают довести эпизод до конца. Поэтому на рис.~10.3 даже первый
    эпизод Sarsa (длиной более 1000 шагов) завершается --- но именно потому, что
    политика улучшается по ходу; переносить эту завершаемость на MC с его
    замороженной политикой нельзя.

    Наконец, даже когда эпизоды завершаются, MC обновляется лишь раз за эпизод по
    полному доходу через тысячи шагов: сигнал размыт, оценки шумные, обучение
    крайне медленное. Бутстрэппинг $n$-шагового Sarsa даёт информативные
    обновления на каждом шаге --- поэтому лучший результат приходится на
    промежуточные $n$ (рис.~10.4), а MC ($n\to\infty$) оказывается наихудшим
    краем. Ключевое преимущество, которого MC лишён, --- возможность учиться внутри ещё не завершённого эпизода.
\end{solution}

\subsection{Упражнение 10.2(с. 294)}
Напишите псевдокод полуградиентного одношагового метода Expected Sarsa для
задачи управления.

\begin{solution}
    Отличие от одношагового Sarsa в цели обновления: вместо $\gamma\,\hat
    q(S',A')$ берётся матожидание по политике $\gamma\sum_a \pi(a\mid S')\hat
    q(S',a,\mathbf w)$. Поскольку конкретное $A'$ в цель больше не входит,
    заранее выбирать и хранить следующее действие не нужно --- действие
    выбирается в начале каждого шага по текущему $\hat q$. В задаче управления
    $\pi$ --- та же $\varepsilon$-жадная по $\hat q(S',\cdot)$ политика, которой
    следует агент (вес $1-\varepsilon+\varepsilon/|\mathcal A|$ на жадном действии
    и $\varepsilon/|\mathcal A|$ на каждом из остальных).

    \begin{algorithm}[H]
        \DontPrintSemicolon
        \SetKwInOut{Input}{Вход}
        \SetKwInOut{Param}{Параметры}
        \Input{дифференцируемая параметризация $\hat q:\mathcal S\times\mathcal A\times\mathbb R^d\to\mathbb R$}
        \Param{размер шага $\alpha>0$, небольшое $\varepsilon>0$}
        Инициализировать веса $\mathbf w\in\mathbb R^d$ произвольно (например, $\mathbf w=\mathbf 0$)\;
        \For{\textnormal{каждого эпизода}}{
            $S\leftarrow$ начальное состояние эпизода\;
            \While{$S$ \textnormal{не заключительное}}{
                Выбрать $A$ как $\varepsilon$-жадное относительно $\hat q(S,\cdot,\mathbf w)$\;
                Предпринять действие $A$; наблюдать $R,S'$\;
                \eIf{$S'$ \textnormal{заключительное}}{
                    $\delta\leftarrow R-\hat q(S,A,\mathbf w)$\;
                }{
                    $\delta\leftarrow R+\gamma\sum_a \pi(a\mid S')\,\hat q(S',a,\mathbf w)-\hat q(S,A,\mathbf w)$\;
                }
                $\mathbf w\leftarrow \mathbf w+\alpha\,\delta\,\nabla\hat q(S,A,\mathbf w)$\;
                $S\leftarrow S'$\;
            }
        }
    \end{algorithm}

    Здесь $\pi(\cdot\mid S')$ --- $\varepsilon$-жадная политика по $\hat
    q(S',\cdot,\mathbf w)$. Если вместо неё взять жадную политику, сумма
    вырождается в $\max_a \hat q(S',a,\mathbf w)$ --- это уже off-policy
    Q-обучение, другой алгоритм.
\end{solution}

\subsection{Упражнение 10.3(с. 294)}
Почему стандартная ошибка для результатов, показанных на рис.~10.4, выше при больших $n$, чем при малых?
\begin{figure}[H]
    \centering
    \includegraphics[width=1\linewidth]{fig_10_4.png}
    \caption{Рис. 10.4. Влияние $\alpha$ и $n$ на качество $n$-шагового полуградиентного Sarsa и аппроксимации функций посредством плиточного кодирования в задаче о машине на
    горе. Как обычно, наилучшее качество достигается при промежуточном уровне бутстрэппинга ($n=4$). Результаты приведены для избранных значений $\alpha$ в логарифмическом масштабе, а затем соединены прямыми линиями. Стандартная ошибка варьировалась от $0.5$ (меньше ширины линии) для $n=1$ до примерно $4$ для $n=16$, поэтому все основные эффекты статистически значимы.}
\end{figure}
\begin{solution}
    Стандартная ошибка здесь --- это $\sigma/\sqrt{N}$, разброс измеренной кривой
    между прогонами ($N=100$ для всех $n$, поэтому различие идёт от $\sigma$ ---
    дисперсии результата между отдельными прогонами).

    Источник --- дисперсия цели обновления. Цель $n$-шагового метода
    \begin{equation*}
        G_{t:t+n} = R_{t+1} + \gamma R_{t+2} + \dots + \gamma^{n-1} R_{t+n}
        + \gamma^n \hat q(S_{t+n}, A_{t+n}, \mathbf w)
    \end{equation*}
    при $n=16$ включает сумму шестнадцати реальных наград, набранных вдоль
    траектории, которая ветвится из-за случайного старта и $\varepsilon$-жадности.
    Чем длиннее этот случайный хвост до бутстрэп-члена, тем выше дисперсия самой
    цели: много случайных шагов накапливают разброс. При $n=1$ цель опирается на
    одну награду плюс бутстрэп и почти детерминирована. (Дисконтирование $\gamma^k$
    при этом гасит вклад далёких наград и работает на снижение дисперсии хвоста ---
    то есть ограничивает рост ошибки, а не вызывает его.)

    Высокодисперсная цель означает, что траектории весов в разных прогонах
    расходятся сильнее, поэтому «число шагов на эпизод» гуляет от прогона к прогону
    шире, и при том же $N$ стандартная ошибка $\sigma/\sqrt{N}$ выше.

    Усиливает эффект и то, что рис.~10.4 усредняет по первым 50 эпизодам --- по
    самой ранней, нестабильной фазе. Там политика и $\hat q$ ещё далеки от
    устойчивых, и длинные $n$-шаговые цели, вычисленные по ещё плохой $\hat q$,
    особенно сильно швыряют результат от прогона к прогону. Поэтому межпрогонная
    дисперсия (а с ней и стандартная ошибка) растёт с $n$ заметнее всего именно на
    этом раннем участке.
\end{solution}

\subsection{Упражнение 10.4(с. 296)}
Напишите псевдокод дифференциальной версии полуградиентного Q-обучения.

\begin{solution}
    Отличие от дифференциального полуградиентного Sarsa в цели: вместо
    $\hat q(S',A')$ берётся $\max_a \hat q(S',a)$ (off-policy, целевая политика ---
    жадная). Поэтому конкретное $A'$ заранее выбирать не нужно --- действие
    выбирается в начале каждого шага по текущему $\hat q$. Тот же $\delta$
    обновляет и веса, и оценку среднего вознаграждения $\bar R$.

    \begin{algorithm}[H]
        \DontPrintSemicolon
        \SetKwInOut{Input}{Вход}
        \SetKwInOut{Param}{Параметры}
        \Input{дифференцируемая параметризация $\hat q:\mathcal S\times\mathcal A\times\mathbb R^d\to\mathbb R$}
        \Param{размеры шагов $\alpha,\beta>0$}
        Инициализировать веса $\mathbf w\in\mathbb R^d$ произвольно (например, $\mathbf w=\mathbf 0$)\;
        Инициализировать оценку среднего вознаграждения $\bar R\in\mathbb R$ произвольно (например, $\bar R=0$)\;
        Инициализировать состояние $S$\;
        \For{\textnormal{каждого шага}}{
            Выбрать $A$ как $\varepsilon$-жадное относительно $\hat q(S,\cdot,\mathbf w)$\;
            Предпринять действие $A$; наблюдать $R,S'$\;
            $\delta\leftarrow R-\bar R+\max_a \hat q(S',a,\mathbf w)-\hat q(S,A,\mathbf w)$\;
            $\bar R\leftarrow \bar R+\beta\,\delta$\;
            $\mathbf w\leftarrow \mathbf w+\alpha\,\delta\,\nabla\hat q(S,A,\mathbf w)$\;
            $S\leftarrow S'$\;
        }
    \end{algorithm}

    Здесь $\bar R$ оценивает средний reward целевой (жадной) политики, поскольку
    $\delta$ построен вокруг $\max_a \hat q(S',a,\mathbf w)$. Это согласовано с
    обновлением весов, в которое входит та же $\bar R$.
\end{solution}

\subsection{Упражнение 10.5(с. 296)}
Какие уравнения (кроме 10.10) необходимы для определения дифференциальной версии алгоритма TD(0)?
\begin{equation}
    \delta_t \doteq R_{t+1} - \bar R_{t+1} + \hat v(S_{t+1}, \mathbf w_t) - \hat v(S_t, \mathbf w_t),
    \tag{10.10}
\end{equation}
\begin{solution}
    Дана дифференциальная TD-ошибка в $v$-форме (10.10):
    \begin{equation*}
        \delta_t = R_{t+1} - \bar R_{t+1} + \hat v(S_{t+1},\mathbf w_t) - \hat v(S_t,\mathbf w_t).
    \end{equation*}
    Кроме неё, для полного определения дифференциального TD(0) нужны ещё два
    уравнения обновления.

    Обновление весов (аналог (10.12), но с градиентом $\hat v$):
    \begin{equation*}
        \mathbf w_{t+1} = \mathbf w_t + \alpha\,\delta_t\,\nabla\hat v(S_t,\mathbf w_t).
    \end{equation*}

    Обновление оценки среднего вознаграждения:
    \begin{equation*}
        \bar R_{t+1} = \bar R_t + \beta\,\delta_t.
    \end{equation*}

    Дополнительно нужно оговорить две вещи. Во-первых, $\delta$ берётся именно в
    $v$-форме (10.10), а не в $q$-форме (10.11) --- то есть оценивается $v_\pi$.
    Во-вторых, это задача \emph{предсказания}: раз мы учим $\hat v$, а не $\hat q$,
    выбор действий по ней невозможен, поэтому политика $\pi$ фиксирована и задана
    извне, а действия просто берутся по ней (никаких правил выбора действия не
    требуется).

    Здесь обновление $\bar R$ через $\delta$ несмещённо, в отличие от
    дифференциального Q-обучения: TD(0) on-policy, $\delta$ построен вокруг $\hat
    v(S_{t+1})$ для той же политики $\pi$, по которой собираются данные, поэтому
    неподвижная точка $\mathbb E[\delta]=0$ даёт ровно $\bar R = r(\pi)$.
\end{solution}

\subsection{Упражнение 10.6(с. 297)}
Рассмотрим марковский процесс вознаграждения, состоящий из трёх состояний $A$, $B$, $C$; переходы производятся детерминированно по кругу. Вознаграждение $+1$ начисляется по прибытии в состояние $A$, в остальных случаях оно равно $0$. Каковы дифференциальные ценности всех трёх состояний?

\begin{solution}
    Цикл $A\to B\to C\to A$ детерминирован, поэтому $\mu(s)=\tfrac13$ для всех
    состояний, а среднее вознаграждение $r(\pi)=\tfrac13$ (одна награда $+1$ на
    три перехода).

    Дифференциальные ценности удовлетворяют уравнению Беллмана без обесценивания
    $v(s) = \big(r_{s\to s'}-r(\pi)\big)+v(s')$, где $+1$ начисляется на переходе,
    прибывающем в $A$ (то есть на $C\to A$):
    \begin{align*}
        v(A) &= \big(0-\tfrac13\big)+v(B),\\
        v(B) &= \big(0-\tfrac13\big)+v(C),\\
        v(C) &= \big(1-\tfrac13\big)+v(A).
    \end{align*}

    Каждое уравнение фиксирует лишь разность соседних ценностей, а их сумма даёт
    тождество $0=0$, поэтому независимых уравнений два на три неизвестных: система
    вырождена, и решение определено с точностью до общей аддитивной константы.
    Это и есть свойство дифференциальных ценностей --- значимы только их разности,
    а общий уровень произволен. Зафиксируем его нормировкой $\sum_s v(s)=0$ и
    получаем
    \begin{equation*}
        v(A) = -\tfrac13, \qquad v(B) = 0, \qquad v(C) = \tfrac13.
    \end{equation*}

    Смысл: ценность тем выше, чем скорее из состояния попадаешь в $A$ (где ждёт
    награда). Ближе всего к $A$ стоит $C$ (один шаг) --- у него ценность
    максимальна; дальше всего $A$ (целый цикл до следующей награды) --- минимальна.
    При другой нормировке (например, $v(A)=0$) получились бы те же разности,
    сдвинутые на константу: $v(A)=0,\ v(B)=\tfrac13,\ v(C)=\tfrac23$.
\end{solution}

\subsection{Упражнение 10.7(с. 298)}
Предположим, что имеется МППР, который при любой стратегии порождает
детерминированную бесконечную последовательность вознаграждений $+1, 0, +1, 0,
\dots$. Технически это недопустимо, потому что нарушается свойство эргодичности
--- не существует стационарного предельного распределения $\mu_\pi$ и предела
(10.7). Тем не менее среднее вознаграждение (10.6) определено корректно. Чему оно
равно? Теперь рассмотрим два состояния этого МППР. При старте из состояния $A$
последовательность вознаграждений точно такая, как описано выше, т.\,е. начинается
с $+1$, а при старте из $B$ последовательность вознаграждений начинается с $0$ и
продолжается так: $+1, 0, +1, 0, \dots$. В этом случае дифференциальный доход
(10.9) определён некорректно, потому что предела не существует. Чтобы исправить
ситуацию, можно было бы определить ценность состояния по-другому:
\begin{equation}
    v_\pi(s) \doteq \lim_{\gamma\to1}\lim_{h\to\infty}
    \sum_{t=0}^{h} \gamma^t\big(\mathbb E_\pi[R_{t+1}\mid S_0=s]-r(\pi)\big).
    \tag{10.13}
\end{equation}
Каковы ценности состояний $A$ и $B$ при таком определении?
\begin{solution}
    Среднее вознаграждение. По (10.6) это предел среднего по первым $h$ шагам.
    Сумма первых $h$ наград последовательности $+1,0,+1,0,\dots$ равна числу
    единиц среди них, то есть $h/2$ с точностью до $\pm\tfrac12$. Деля на $h$:
    \begin{equation*}
        r(\pi) = \lim_{h\to\infty}\frac1h\sum_{t=1}^h R_t
        = \lim_{h\to\infty}\Big(\tfrac12 \pm \tfrac{1}{2h}\Big) = \tfrac12.
    \end{equation*}
    Сдвиг фазы между стартами из $A$ и $B$ меняет частичную сумму лишь на единицу
    --- фиксированно, тогда как делим на растущее $h$, --- поэтому в пределе он
    вымывается, и $r(\pi)=\tfrac12$ для обоих стартов. При этом $r(\pi)$ существует,
    хотя $\mu_\pi$ и предел (10.7) --- нет: усреднение по времени работает там, где
    усреднение по стационарному распределению неприменимо.

    Ценности через (10.13). Из $A$ отклонения наград от $r(\pi)$ равны
    $\tfrac12(-1)^t$, $t=0,1,2,\dots$. Подставляя в (10.13) и пользуясь
    $\sum_{t=0}^\infty(-\gamma)^t = \tfrac{1}{1+\gamma}$ при $\gamma<1$:
    \begin{equation*}
        v(A) = \lim_{\gamma\to1}\sum_{t=0}^\infty \gamma^t\cdot\tfrac12(-1)^t
        = \lim_{\gamma\to1}\frac{1}{2(1+\gamma)} = \frac14.
    \end{equation*}
    Из $B$ отклонения равны $\tfrac12(-1)^{t+1}$, что даёт противоположный знак:
    \begin{equation*}
        v(B) = \lim_{\gamma\to1}\Big(-\frac{1}{2(1+\gamma)}\Big) = -\frac14.
    \end{equation*}

    Именно ради этого в (10.13) предел по $h$ берётся при $\gamma<1$ (ряд
    сходится), и лишь затем $\gamma\to1$: знакочередующийся ряд $+\tfrac12-\tfrac12
    +\tfrac12-\dots$ суммируется по Абелю к середине между частичными суммами
    $\tfrac12$ и $0$, то есть к $\tfrac14$. Итог: $v(A)=+\tfrac14$,
    $v(B)=-\tfrac14$ --- из $A$ награда приходит сразу, поэтому его ценность выше.
\end{solution}

\subsection{Упражнение 10.8(с. 298)}
Псевдокод во врезке на стр.~296 обновляет $\bar R_{t+1}$, используя в качестве
ошибки $\delta_t$, а не просто $R_{t+1}-\bar R_{t+1}$. Годятся оба определения
ошибки, но $\delta_t$ даёт лучшие результаты. Чтобы понять, почему, рассмотрим
круговой МППР с тремя состояниями из упражнения 10.6. Оценка среднего
вознаграждения должна стремиться к своему истинному значению $1/3$. Предположим,
что она уже принимала это значение и там и застряла. Какой тогда была бы
последовательность $R_t-\bar R_t$? А какой была бы последовательность $\delta_t$
(вычисленная по формуле (10.10))? Какая последовательность давала бы более
устойчивую оценку среднего вознаграждения, если бы оценке было разрешено
изменяться в ответ на ошибки? Объясните свой ответ.

\begin{equation}
    \delta_t \doteq R_{t+1} - \bar R_{t+1} + \hat v(S_{t+1}, \mathbf w_t) - \hat v(S_t, \mathbf w_t).
    \tag{10.10}
\end{equation}

\begin{solution}
    Пусть $\bar R$ застряла в истинном значении $\tfrac13$. Используем истинные
    дифференциальные ценности из упражнения 10.6: $v(A)=-\tfrac13$, $v(B)=0$,
    $v(C)=\tfrac13$ (награда $+1$ только на входе в $A$, то есть на переходе
    $C\to A$).

    Последовательность $R_t-\bar R_t$. Награды по циклу $A\to B\to C\to A$ суть
    $0,0,1$, поэтому
    \begin{equation*}
        R_t - \bar R_t:\quad -\tfrac13,\ -\tfrac13,\ +\tfrac23,\ -\tfrac13,\ -\tfrac13,\ +\tfrac23,\ \dots
    \end{equation*}
    Она колеблется внутри цикла, хотя $\bar R$ уже верна.

    Последовательность $\delta_t$. По (10.10) $\delta_t = R_{t+1}-\bar R +
    \big(v(S_{t+1})-v(S_t)\big)$. На каждом переходе:
    \begin{align*}
        A\to B:&\quad 0-\tfrac13+(0-(-\tfrac13)) = 0,\\
        B\to C:&\quad 0-\tfrac13+(\tfrac13-0) = 0,\\
        C\to A:&\quad 1-\tfrac13+(-\tfrac13-\tfrac13) = 0.
    \end{align*}
    То есть $\delta_t \equiv 0$ --- тождественно, на каждом шаге, а не только в
    среднем. Бутстрап-член $v(S_{t+1})-v(S_t)$ поглощает внутрицикловые
    колебания награды: всплеск $+1$ «объясняется» входом в более ценное состояние,
    а не списывается на ошибку среднего.

    Какая оценка устойчивее. Пусть $\bar R$ разрешено двигаться: $\bar R \leftarrow
    \bar R + \beta\cdot(\text{ошибка})$. С ошибкой $R_t-\bar R_t$ толчки
    $-\tfrac13, -\tfrac13, +\tfrac23$ ненулевые даже в правильной точке, поэтому
    $\bar R$ постоянно дёргается вокруг $\tfrac13$ --- лишняя дисперсия на ровном
    месте. С ошибкой $\delta_t\equiv 0$ в правильной точке двигать $\bar R$ нечем,
    и она там остаётся.

    Поэтому $\delta_t$ даёт более устойчивую оценку: в неподвижной точке она равна
    нулю и не расшатывает уже верную $\bar R$, тогда как сырая невязка $R_t-\bar
    R_t$ расшатывает. Бутстрап-член снимает предсказуемые внутрицикловые
    колебания, оставляя в $\delta_t$ только настоящую ошибку среднего.
\end{solution}

\subsection{Упражнение 10.9(с. 302)}
В дифференциальном полуградиентном $n$-шаговом алгоритме Sarsa размер шага для среднего вознаграждения $\beta$ должен быть очень малым, чтобы $\bar R$ была хорошей оценкой в долгосрочной перспективе. К сожалению, тогда $\bar R$ для многих шагов будет смещена на своё начальное значение, что делает обучение неэффективным. Альтернативно можно было бы использовать выборочное среднее, но из-за нестационарности (стратегия дрейфует) оно не годится. Размер шага для среднего вознаграждения --- идеальное место для приёма с несмещённым постоянным шагом из упражнения 2.7. Опишите конкретные изменения в алгоритме.

\begin{solution}
    Заменяем постоянный $\beta$ на динамический шаг $\beta_n=\alpha/\bar o_n$ из
    (2.8)--(2.9), и только для оценки среднего вознаграждения $\bar R$; шаг
    $\alpha$ для весов $\mathbf w$ остаётся обычной константой. Конкретно:

    \begin{itemize}
        \item Во вход вместо $\beta$ берём базовый $\alpha>0$ и добавляем
            аккумулятор $\bar o$, инициализированный нулём ($\bar o_0\doteq 0$).
        \item В ветке $\tau\ge 0$, перед обновлением $\bar R$, продвигаем $\bar o$
            и пересчитываем шаг:
            \begin{equation*}
                \bar o \leftarrow \bar o + \alpha(1-\bar o), \qquad
                \beta \leftarrow \alpha/\bar o.
            \end{equation*}
        \item Обновляем среднее этим $\beta$: $\ \bar R \leftarrow \bar R + \beta\,\delta$.
    \end{itemize}

    Порядок важен: $\bar o$ обновляется \emph{до} вычисления $\beta$, так как
    (2.9) использует уже продвинутое $\bar o_n$. Иначе на первом шаге было бы
    деление на $\bar o_0=0$. После первого шага $\bar o_1=\alpha$, значит
    $\beta_1=\alpha/\alpha=1$: первое обновление $\bar R$ идёт с полным весом и
    полностью стирает произвольное начальное $\bar R_0$ (по анализу упражнения 2.7
    множитель $1-\beta_1=0$ обнуляет всё произведение $\prod(1-\beta_i)$).

    Фрагмент обновления в ветке $\tau\ge 0$ принимает вид:
    \begin{align*}
        \delta &\leftarrow \sum_{i=\tau+1}^{\tau+n}\big(R_i-\bar R\big)
                 + \hat q(S_{\tau+n},A_{\tau+n},\mathbf w) - \hat q(S_\tau,A_\tau,\mathbf w)\\
        \bar o &\leftarrow \bar o + \alpha(1-\bar o)\\
        \beta  &\leftarrow \alpha/\bar o\\
        \bar R &\leftarrow \bar R + \beta\,\delta\\
        \mathbf w &\leftarrow \mathbf w + \alpha\,\delta\,\nabla\hat q(S_\tau,A_\tau,\mathbf w)
    \end{align*}

    Так $\bar R$ получает несмещённость (начальное значение стирается на первом
    шаге) и сохраняет восстановительную способность постоянного шага, нужную при
    дрейфе стратегии, --- в отличие и от постоянного $\beta$ (вечный след $\bar
    R_0$), и от выборочного среднего (не справляется с нестационарностью).
\end{solution}
